\section{No five-term equation of width at most 65,536}
The \emph{width} of a polynomial is the difference between its largest and smallest exponents with nonzero coefficients. For a nonzero singular modulus we may divide the equation by its lowest power of $X$ and write $0=e_0<e_1<e_2<e_3<e_4$.

\begin{corollary}\label{cor:width}
Every singular modulus of degree at least five requires at least six terms in a rational equation of width at most $65{,}536$.
\end{corollary}
\begin{proof}
By Theorems~\ref{th:four} and~\ref{th:five} only five discriminants, and exactly five terms, require a test. At each discriminant the table gives four roots $u_1,\ldots,u_4$ with $v_p(u_i/r-1)>3/(p-1)$ and an additional root $w$ with $v_p(w)=v>0$. The checker verifies both assertions from the class polynomial. For the first, apply the least-minimizer test after translation; for the second, verify that the minimum of $v_p(A_i)+iv$ occurs at two or more indices.
\begin{center}\small
\begin{tabular}{@{}rrrrrr@{}}\toprule
$\Delta$ & $p$ & $r$ & $v$ & Pairs tested & Second-prime rechecks \\ \midrule
$-87$  & 23 & 19 & 3   & 2,977 & 2,983 \\
$-131$ & 17 &  8 & 3   & 4,094 & 3,994 \\
$-348$ & 29 &  2 & $3/2$ & 4,676 & 4,829 \\
$-432$ & 23 & 19 & 3   & 2,977 & 2,980 \\
$-675$ & 17 &  8 & 3   & 4,094 & 4,191 \\ \bottomrule
\end{tabular}
\end{center}
The four-row cofactor calculation from the preceding section still holds. Indeed every $x_\ell=\log(u_\ell/r)$ has valuation strictly greater than $3/(p-1)$. For $0\le i\le3$ and $m>i$, the factorial estimates proving (1) bound $v_p(m!/i!)/(m-i)$ by $1/(p-3)\le 3/(p-1)$. All nonleading Schur terms have strictly positive valuation after the two Vandermonde factors are removed. Expanding along $w$ now gives a pair $1\le j<i\le4$ with
\[
          p^{\lceil ve_j/3\rceil}\mid e_i-e_j .
\]
The larger exponent is indeed $e_i$, since $p^{\lceil v m/3\rceil}>m$ for every positive integer $m$ in these five cases.

For each admissible pair of exponents in $[1,65536]$, reduce all powers modulo the monic class polynomial over the two prime fields from the preceding section. Quotient out the constant column. For degree six project the remaining five coefficients $(a_1,\ldots,a_5)$ to $(a_1+a_5,\ldots,a_4+a_5)$; for degree five keep all four. A full-rank image proves independence. In the two-dimensional quotient by the fixed pair, distinct projective directions for the other two columns prove full rank at the first prime. All coincident directions are rechecked at the second prime. Every recorded recheck gives full rank. This verifies the claimed width for all five possible exceptions.
\end{proof}
